\begin{table*}[htbp]
\centering
\small
\begin{tabular}{lrrrrrrr}
\toprule
& & \multicolumn{3}{c}{\textbf{Global view}} & \multicolumn{3}{c}{\textbf{Within-stratum view}} \\
\cmidrule(lr){3-5} \cmidrule(lr){6-8}
\textbf{Train samples} & \textbf{\# Langs} & \unilid & calib.\ \unilid & \fasttext & \unilid & calib.\ \unilid & \fasttext \\
\midrule
$<$500 & 56 & 0.515 & 0.780 & 0.750 & 0.871 & 0.827 & 0.915 \\
500--1k & 40 & 0.628 & 0.892 & 0.861 & 0.975 & 0.955 & 0.964 \\
1k--12k & 458 & 0.891 & 0.945 & 0.941 & 0.990 & 0.987 & 0.979 \\
12k--18k & 526 & 0.979 & 0.985 & 0.971 & 0.997 & 0.997 & 0.986 \\
18k--35k & 398 & 0.963 & 0.965 & 0.953 & 0.992 & 0.992 & 0.981 \\
35k+ & 462 & 0.958 & 0.957 & 0.941 & 0.958 & 0.958 & 0.942 \\
\bottomrule
\end{tabular}
\caption{Mean per-language F1 by training-resource tier under both metric
views, on the 45{,}377{,}279 scored test lines. Global view: every false
positive counts against the predicted language. Within-stratum view: test
examples are restricted to true labels inside the tier, so false positives
arriving from other tiers are excluded (the view used by
\cref{tab:resource-tier}). The two views rank the methods oppositely in the
smallest tiers: calibrated \unilid removes false positives \camrev{\emph{into}} small
languages, which only the global view counts, and re-examination moves some
of the smallest languages' own low-margin lines away, which only the
within-stratum view penalizes.}
\label{tab:calibrated_views}
\end{table*}
